problem:  
Let \(G\) be a connected, simply connected **real semisimple Lie group** with Lie algebra \(\mathfrak{g}\).  
Given a **left‑invariant Lorentzian metric** \(g\) on \(G\) corresponding to a nondegenerate symmetric bilinear form \(B\) of signature \((1,\dim G-1)\) on \(\mathfrak g\), the **Euler–Arnold equation** governing its geodesics reads  
\[
\dot{u} = -\,\operatorname{ad}(u)^{\top}u,\qquad u(t)\in\mathfrak g,
\]
where \(\operatorname{ad}(u)^{\top}\) is the adjoint of \(\operatorname{ad}(u)\) with respect to \(B\).  
Each flow line \(u(t)\) corresponds to a (possibly incomplete) geodesic on \((G,g)\).  

**Problem:**  
For semisimple \(\mathfrak g\), the Riemannian case (when \(B\) is positive definite) is always geodesically complete, but in the Lorentzian case incompleteness can occur even for compact \(G\).  

The open question is:  
Can the **geodesic completeness** of the Lorentzian left‑invariant metric \((G,g)\) be **detected purely from the adjoint‑invariant polynomial structure** on \(\mathfrak g\)?  

More concretely, does there exist a description—perhaps via the classical ring of invariant polynomials \(\mathbb R[\mathfrak g]^{\operatorname{Ad}(G)}\)—of conditions on \(B\) ensuring that every trajectory of  
\[
X_B(u) = -\,\operatorname{ad}(u)^{\top}u
\]
is complete?  Equivalently, is completeness equivalent to the existence of a **complete Hamiltonian vector field** on \(\mathfrak g^{*}\) whose Hamiltonian and symplectic form are both induced from \(B\) and the Lie–Poisson structure?  

An even sharper formulation asks:  
> For which forms \(B\) on a real semisimple \(\mathfrak g\) is the associated vector field \(X_B\) complete, and can these be characterized in terms of the behavior of \(B\) on the coadjoint orbits defined by the invariant polynomials?

Why is it a "good" problem:  
This problem isolates a precise and foundational issue at the interface between **Lorentzian geometry**, **Lie theory**, and **Hamiltonian dynamics**. Rather than seeking purely analytic or case‑by‑case criteria, it asks whether the **algebraic invariant structure of a semisimple Lie algebra**—embodied in its Casimir polynomials and coadjoint orbits—is sufficient to determine a **global analytic property** (geodesic completeness) of a Lorentzian metric. It thus bridges **invariant theory** and **global analysis** in a way not yet understood. Conceptually simple (“Is completeness algebraically visible?”), it targets the deepest layer of the Lorentzian incompleteness phenomenon—whether it stems from algebraic instability or from analytic blow‑up—offering a potential doorway to unifying algebraic and geometric completeness theories on Lie groups.